\begin{remarks}
\label{II.4.14}
Criterion (iii) of Theorem~\ref{II.4.10} deserves to be called
the \emph{Jacobian criterion of smoothness}.
It allows one to recognize, theoretically, whether a given
$S$-prescheme $Y$ is smooth over~$S$ at a point $x$
of~$Y$, since there always exists a neighborhood of~$Y$ isomorphic to
a subprescheme of a smooth $S$-prescheme $X$, for
example $X=S[t_1,\dots,t_n]$. It is moreover for
$X=S[t_1,\dots,t_n]$, $S=\Spec(A)$, that one usually states the
Jacobian criterion
% original p. 44
(of course, in the classical case envisaged by Zariski, $A$ was
a field). One leaves to the reader to give the statement relative
to the datum of an ideal $J$ of~$A[t_1,\dots,t_n]$ and of a
prime ideal containing it, to which one is thus led. Let us note
that it indeed seems at the present time (and especially since Nagata
has succeeded in generalizing by
non-differential methods Zariski's theorem stating that
the set of regular points of an algebraic scheme is
open) that the Jacobian criterion is scarcely of interest
except in the form in which we give it here (\ie using
exclusively \emph{relative} differentials and not
\emph{absolute} differentials, \ie relative to the ring of
absolute constants $\ZZ$). As so often, the consideration
of differentials is more convenient here than that of
derivations. Let us note finally that if $Y$ is smooth over~$S$ at $x$, of
relative dimension $n$, then there exists an open neighborhood of~$x$
in~$Y$ isomorphic to a subprescheme of
$X=S[t_1,\dots,t_n]$ with $n=m+1$, as follows from the
definition and
from~(I~\SourceRef{I.7.6}{7.6}).

Let $A$ be a noetherian ring, $x_i$ ($1 \leq i \leq n$)
elements of~$A$, $J$ the ideal generated by the
$x_i$. One says that the $x_i$ form a \emph{regular
system of generators}
of~$J$ if the canonical surjective homomorphism
\[
(A/J)[t_1,\dots,t_n] \to \gr^J(A)
\]
defined by the $x_i$ (where the second member denotes
the graded ring associated to $A$ filtered by the
powers of~$J$) is an \emph{isomorphism}. This condition
also means that
\begin{enumerate}
\item[(i)] The canonical surjective homomorphism
\[
S_{A/J} (J/J^2) \to \gr^J(A)
\]
(where the first member denotes the symmetric algebra of the
$A/J$-module $J/J^2$) is an isomorphism, and
\item[(ii)] $J/J^2$ is free and admits as basis the classes of the $x_i
\mod J^2$.
\end{enumerate}
In this form, one sees that if $J\neq A$, the $x_i$ form a
\emph{minimal system of generators} of~$J$, and that
\emph{every other minimal system of generators} of~$J$
\emph{is a regular system of generators}
(N.B.\ ``minimal'' is taken in the strict sense: minimum number
of elements, which is equivalent to the sense: minimal for
inclusion, only if $A$ is local); on the other hand, if $J=A$, every
system of generators of~$J$ is regular.

The
% original p. 45
condition of regularity of a system of
generators of an ideal is stable under localization with respect to an
arbitrary multiplicatively closed set, and on the other hand one sees
at once that for $(x_i)$ to be a minimal system of
generators of~$J$, it already suffices that for every
\emph{maximal ideal $\goth{m}$ containing $J$}, the $x_i$ define a
regular system of generators of~$JA_{\goth{m}}$ in
$A_{\goth{m}}$. This therefore reduces us to the case where $A$ is a
local ring with maximal ideal $\goth{m}$, and where the $x_i$ are
in $\goth{m}$. \emph{Then the $x_i$ form a
regular system of generators of~$J$ if and only if they
form an $A$-sequence in the sense of Serre}\footnote{We shall
now rather say ``$A$-regular sequence'', \cf EGA
$0_{\textup{IV}}$~15.1.7 and 15.1.11.}, \ie if for every $i$ such that $1
\leq i \leq n$, $x_i$ is not a divisor of~$0$ in
$A/(x_1,\dots,x_{i-1})A$.

Finally, in the case where $A$ is an algebra over a ring $B$,
and where $A/J$ is isomorphic as a $B$-algebra to $B$ (so
that $J$ is the kernel of a homomorphism of $B$-algebras
$A\to B$), then the $x_i$ form a regular system of
generators of~$J$ if and only if the canonical
homomorphism
\[
B[[t_1,\dots,t_n]]\to \hat{A}
\]
defined by the $x_i$ (where the second member denotes the
separated completion $\varprojlim A/J^{n+1}$ of~$A$ for the
topology defined by the powers of~$J$) is an
\emph{isomorphism} (it is in any case \emph{surjective}).

All these facts are well known, and doubtless appear in Serre's
course on commutative algebra written up by Gabriel, up
to minor details. (Where one finds $N$ other
characterizations of $A$-sequences, in the case where $A$ is a
local ring).

Let $J$ be an ideal in a noetherian ring $A$. We shall say that
$J$ is a \emph{regular ideal}
if for every prime ideal $\goth{p}$ of~$A$, $JA_{\goth{p}}$ admits
a regular system of generators. It is
obviously sufficient to verify this for $\goth{p} \supset J$, and one may
moreover restrict oneself to maximal $\goth{p}$. More generally,
let $\cal{J}$ be an ideal on a locally
noetherian prescheme $X$, one says that $\cal{J}$ is a \emph{regular
ideal} if for every $x\in X$, $\cal{J}_x$ is an ideal of
$\cal{O}_x$ which admits a regular system of
generators. This is equivalent to the conjunction of the two
following conditions:
\begin{enumerate}
\item[(a)] The canonical surjective homomorphism
\[
S_{\cal{O}_X/\cal{J}}(\cal{J}/\cal{J}^2) \to \gr^{\cal{J}}(\cal{O}_X)
\]
is an isomorphism and
\item[(b)] The sheaf of~$\cal{O}_X/\cal{J}$-Modules
$\cal{J}/\cal{J}^2$ is locally free.
\end{enumerate}

One
% original p. 46
then also says that the subprescheme $Y$ of~$X$ defined
by $\cal{J}$ (hence such that $\cal{O}_Y$ extended by $0$ is
isomorphic to $\cal{O}_X/\cal{J}$) is \emph{regularly
immersed}
in~$X$, and one defines in the same way (in an evident
manner) the notion of morphism of \emph{regular
immersion},
(\resp \emph{regular at a point $x$}),
immersion morphism $Y\to X$ identifying $Y$ (\resp a suitable neighborhood of
$x$) with a closed subprescheme regularly
immersed in an open subset of~$X$. (One must not say:
regular subprescheme, for that would mean that the
local rings of~$Y$ are regular). Finally, sections $x_i$
of~$\cal{J}$ are called a \emph{regular system of
generators} if for every $x\in X$, the corresponding
elements of~$\cal{O}_x$ form a regular system of
generators of~$\cal{J}_x$ (terminology compatible with that
introduced for generators of an ideal of a
ring). This also means that the canonical surjective homomorphism
\[
\cal{O}_Y[t_1,\dots,t_n] \to \gr^{\cal{J}} (\cal{O}_X)
\]
defined by the $x_i$ is an isomorphism. If one knows in advance
that the ideal $\cal{J}$ is regular, this also means,
simply, that at every point $x$ \emph{of~$Y$}, the $x_i$
define a \emph{basis} of~$\cal{J}/\cal{J}^2$ over~$\cal{O}_{Y,x}$. (N.B.\ this condition is empty if $Y$ is
empty). Thus, for $\cal{J}$ to admit a regular
system of generators, it is necessary and sufficient that $\cal{J}$ be
regular, and that the $\cal{O}_Y$-Module $\cal{J} /\cal{J}^2$ be
globally free (and not only locally free), \ie that
the canonical homomorphism $S_{\cal{O}_Y}(\cal{J}/\cal{J}^2) \to
\gr^{\cal{J}} (\cal{O}_X)$ be surjective, and that the
$\cal{O}_Y$-Module $\cal{J}/\cal{J}^2$ be globally free.

An \emph{augmented ring is said to be regular}
if the ideal of the augmentation is regular. Thus, if $A$ is
a local ring, regarded as augmented in its residue
field $k$, then $A$ is a regular local ring if and
only if it is a regular augmented ring.

(To tell the truth, it seems that it was useless to begin by
doing the preliminary sorites for rings; there is an
advantage in beginning with sheaves at once. If
one wants something in the noetherian case, it is the
definition adopted here---a priori less strict than that by
Serre's $A$-sequences---which seems preferable for the
needs of differential calculus. Of course, to do things properly, one
should also develop at least part of the theory
of smooth morphisms in the non-noetherian setting\footnote{As
is said in the foreword, this has now been done, \cf EGA
IV 17, 18.}, probably
% original p. 47
starting from the Jacobian criterion, so as to obtain if
possible all the essential formal properties of
smooth morphisms and of \'etale morphisms \ie smooth and
quasi-finite; the converses alone appealing to
noetherian hypotheses).
\end{remarks}
